\documentclass[11pt]{article}

\usepackage[T1]{fontenc}
\usepackage[utf8]{inputenc}
\usepackage[english]{babel}
\usepackage{amsmath,amssymb,amsthm}
\usepackage{booktabs}
\usepackage{array}
\usepackage{placeins}
\usepackage{geometry}
\usepackage{hyperref}
\usepackage{url}

\geometry{margin=1in}
\setlength{\emergencystretch}{2em}

\hypersetup{
  colorlinks=true,
  linkcolor=black,
  citecolor=blue,
  urlcolor=blue,
  pdftitle={Independent Exhaustive Verification and a Formula-Derived b-File for OEIS A369431},
  pdfauthor={Carlo Corti}
}

\newtheorem{theorem}{Theorem}
\newtheorem{proposition}[theorem]{Proposition}
\newtheorem{lemma}[theorem]{Lemma}
\newtheorem{corollary}[theorem]{Corollary}
\theoremstyle{definition}
\newtheorem{definition}[theorem]{Definition}
\theoremstyle{remark}
\newtheorem{remark}[theorem]{Remark}
\newcolumntype{L}[1]{>{\raggedright\arraybackslash}p{#1}}

\title{Independent Exhaustive Verification of OEIS A369431 through $n=22$\\
and a Formula-Derived b-File through $n=400$}
\author{Carlo Corti\\
Independent researcher\\
\texttt{carlo.corti@outlook.com}}
\date{July 22, 2026}

\begin{document}

\maketitle

\begin{abstract}
OEIS sequence A369431 counts permutations of $[n]$ avoiding the four
classical patterns $1234$, $1324$, $1342$, and $2413$.  An independently
implemented, formula-free maximum-insertion program exhaustively computes
\[
 a(21)=8321204878,
 \qquad
 a(22)=26674199972.
\]
Every avoiding permutation is reached exactly once by repeatedly inserting
the next maximum into an avoiding parent.  Because the parent already avoids
the forbidden set, the incremental predicate need inspect only new occurrences that
contain the inserted maximum.  Exhaustive coverage is organized through the
20 order-$4$ avoiders and their 66 order-$5$ children, while a small fixed
number of operational groups provides restartable human-scale runs without
fragmenting the mathematical computation or its archival record.  The two
reported totals are supported by append-only manifests, final aggregate
tables, small-order direct oracles and replay tests, sanitizer and static
analysis, and external LLM-assisted code reviews.  A subsequent literature
search found that PermPAL had already recorded a verified combinatorial
specification and rational generating function for this exact class.  We use
that prior result, with explicit attribution, to produce an OEIS-style b-file
through $n=400$.  Exact Python recurrence iteration and an independent
PARI/GP formal-series expansion agree for all 401 coefficients; the first
101 also agree with the sequence displayed by PermPAL.  No priority is claimed
for the formula, recurrence, or coefficients.  The contribution is instead a
transparent separation of a finite independent enumeration through $n=22$
from a reproducible formula-derived data extension thereafter.
\end{abstract}

\noindent\textbf{2020 Mathematics Subject Classification.}
Primary 05A05; Secondary 05A15, 11Y55.

\noindent\textbf{Key words and phrases.}
permutation patterns, pattern avoidance, exact enumeration, rational
generating function, integer sequence, OEIS, reproducible computation.

\section{Introduction}
\label{sec:introduction}

We write $[n]=\{1,\ldots,n\}$ and call $n$ the order (or length) of a
permutation.  Let
\[
 B=\{1234,1324,1342,2413\}
\]
and let $\operatorname{Av}_n(B)$ denote the permutations of $[n]$ that avoid
every classical pattern in $B$.  Since the four distinct forbidden patterns
have the same length, they form an antichain, and each is necessary; hence
$B$ is also the basis of this permutation class.  The OEIS entry A369431 is
\begin{equation}
 a(n)=\lvert\operatorname{Av}_n(B)\rvert .
 \label{eq:definition}
\end{equation}
The local OEIS snapshot has offset $0$, with $a(0)=1$ for the empty
permutation, and gives the terms through
\[
 a(20)=2595858574.
\]
It credits Matt Slattery-Holmes with submitting the sequence on January 23,
2024, and Martin Ehrenstein with extending it from $a(13)$ through $a(20)$ on
February 24, 2024~\cite{oeisA369431}.  For example, at $n=4$ precisely the four
members of $B$ are excluded from the 24 permutations, so $a(4)=20$.

Classical pattern avoidance belongs to the established theory of restricted
permutations, with foundational results due to Simion and Schmidt and modern
book-length treatments by B\'ona and Kitaev
\cite{simionSchmidt1985,bona2022,kitaev2011}.  MathWorld provides a concise
online introduction and links pattern-avoidance counts to their OEIS records
\cite{weissteinPermutationPattern}.  These sources supply the general context.

A literature and database search prompted by the first review round located
this exact class in the Permutation Pattern Avoidance Library (PermPAL), where
it is accompanied by an automatically verified combinatorial specification,
a rational generating function, and a counting sequence that includes the two
values computed here~\cite{albertEtAl2026,permpalA369431Class}.  Thus the
novelty claimed here is not priority for those integers, the recurrence, or
the generating function.  The contribution is the independent, formula-free
exhaustive check through $n=22$, its compact audit trail, and a separately
labelled, reproducible b-file calculation from the prior structural result.

The principal difficulty in a direct extension is the number of avoiding
objects themselves.  Enumerating all $n!$ permutations is unnecessary, but a
formula-free generator must still reach every avoiding permutation and must
do so without omission or duplication.  The computation described here uses
the hereditary nature of classical pattern avoidance: it starts from a small
avoiding permutation and inserts successive maxima, pruning an insertion as
soon as it creates one of the four forbidden patterns.

Successive insertion is also central to the insertion-encoding approach to
permutation classes developed by Albert, Linton, and Ru\v{s}kuc
\cite{albertLintonRuskuc2005}.  The present method is deliberately narrower:
it does not construct a regular-language encoding, infer a generating
function, or consume the PermPAL specification.  It uses a fixed
maximum-insertion tree and a basis-specific direct predicate solely for exact
finite enumeration.

The contributions of this paper are:
\begin{enumerate}
  \item a direct correctness argument for a sequence-specific incremental
        maximum-insertion predicate;
  \item a complete and disjoint coverage scheme based on 20 mathematical
        atoms of order $4$, refined into 66 durable internal units of order
        $5$;
  \item an operational grouping and persistence design that retains one
        monolithic production source, one manifest per target, and one final
        aggregate per target;
  \item an artifact-bounded, independently implemented computation of the two
        values stated in Theorem~\ref{thm:main}; and
  \item a formula-derived b-file through $n=400$, cross-checked by two exact
        calculation paths and against all 101 coefficients displayed by
        PermPAL.
\end{enumerate}

The exhaustive C computation does not depend on a recurrence or generating
function.  The later b-file calculation in Section~\ref{sec:recurrence} does;
the two evidential tracks are kept explicit throughout.

\section{Pattern avoidance and maximum insertion}
\label{sec:mathematics}

\begin{definition}
Let $\pi=\pi_1\cdots\pi_n$ be a permutation and let
$\sigma=\sigma_1\cdots\sigma_k$ be a pattern.  The permutation $\pi$
\emph{contains} $\sigma$ if some subsequence
$\pi_{i_1},\ldots,\pi_{i_k}$, with
$i_1<\cdots<i_k$, has the same relative order as $\sigma$.  Otherwise $\pi$
\emph{avoids} $\sigma$.
\end{definition}

Suppose that $\pi\in\operatorname{Av}_m(B)$ and form $q$ by inserting the new
maximum $M=m+1$ at position $p$.  An occurrence of a forbidden pattern that
does not use $M$ would already occur in $\pi$.  Thus a new forbidden
occurrence must contain $M$.  Since $M$ is the largest of the four values, it
corresponds to the entry $4$ in each forbidden pattern.  The four possible
cases are exactly
\begin{align}
1234 &: i<j<k<p, & q_i&<q_j<q_k<M, \label{eq:case1234}\\
1324 &: i<j<k<p, & q_i&<q_k<q_j<M, \label{eq:case1324}\\
1342 &: i<j<p<k, & q_i&<q_k<q_j<M, \label{eq:case1342}\\
2413 &: i<p<j<k, & q_j&<q_i<q_k<M. \label{eq:case2413}
\end{align}
The production predicate is the negation of the existence of any index tuple
satisfying one of \eqref{eq:case1234}--\eqref{eq:case2413}.  It is therefore a
fixed predicate for this basis, not a generic pattern parser.

\begin{proposition}[Correctness of maximum insertion]
\label{prop:insertion}
For every $n\geq 1$, recursively inserting the next maximum in every position
that passes the four tests above generates every member of
$\operatorname{Av}_n(B)$ exactly once and generates no other permutation.
\end{proposition}

\begin{proof}
Proceed by induction on the order.  The initial permutation of order $1$ is
the unique member of $\operatorname{Av}_1(B)$ and avoids $B$.  If an avoiding
parent is extended, an old forbidden occurrence
cannot appear; equations \eqref{eq:case1234}--\eqref{eq:case2413} enumerate
all possible new forbidden occurrences containing the inserted maximum.
Consequently, the predicate accepts exactly the avoiding children.

Conversely, delete the maximum from any $q\in\operatorname{Av}_{m+1}(B)$.
The result is an avoiding permutation $\pi$ of order $m$, because classical
pattern avoidance is hereditary.  Reinserting $m+1$ at its original position
recovers $q$ and passes the incremental test.  Deletion of the maximum is a
unique inverse because the maximum has one uniquely determined position.
Thus the construction is both complete and duplicate-free.  Heredity also
ensures that every intermediate permutation on an insertion path is an
avoider, so the incremental premise is valid at every level.  The recursion
terminates because the order increases by one at every step and stops at the
fixed target $n$.
\end{proof}

This proposition establishes the mathematical enumeration scheme.  The
finite numerical conclusions additionally depend on the faithful C17
implementation and the recorded completed runs, whose validation boundary is
described below.

\section{Coverage atoms, internal units, and operational groups}
\label{sec:coverage}

\subsection{Twenty mathematical coverage atoms}

There are $a(4)=20$ avoiders of order $4$.  They form a natural disjoint
frontier: repeatedly deleting the current maximum from any avoiding
permutation of order at least $4$ ends at exactly one member of
$\operatorname{Av}_4(B)$.  The program stores these 20 roots explicitly in
lexicographic order.  They are \emph{mathematical coverage atoms}; they are
not 20 source files, 20 manifests, or a requirement for 20 operator launches.

Indeed, Proposition~\ref{prop:insertion} shows that deleting the maximum maps
an avoider to its unique avoiding parent.  Iterating this map down to order
$4$ assigns every target-order avoider to exactly one of the 20 roots.
Consequently the 20 rooted subtrees are pairwise disjoint and their union is
$\operatorname{Av}_n(B)$ for every $n\geq4$.

\subsection{The order-5 recovery frontier}

For targets above $4$, each valid insertion of $5$ into the 20 roots defines
one internal unit.  Roots are numbered $0,\ldots,19$ in lexicographic order;
the five insertion positions are numbered $0,\ldots,4$, from before the first
entry through after the last.  Ordered first by root and then by insertion
position, the accepted children are numbered $u=0,\ldots,65$ and are exactly
$\operatorname{Av}_5(B)$; hence they number
\[
 a(5)=66.
\]
By Proposition~\ref{prop:insertion}, their descendant subtrees are complete,
pairwise disjoint, and collectively cover the target order.  The refinement
from 20 atoms to 66 units is useful for recovery: after an interruption, a
completed unit need not be counted again.  It does not change the underlying
mathematical partition.

\subsection{Segmentation without fragmentation}
\label{subsec:groups}

The 66 units are bundled into a small fixed number $G$ of operational groups,
numbered $0,\ldots,G-1$:
$G=4$ for the $n=21$ campaign and $G=8$ for the $n=22$ campaign.  Unit $u$
belongs to group $u\bmod G$, which interleaves neighboring lexicographic units
and reduces the imbalance that would result from assigning contiguous
subtrees.  One command processes one entire operational group and then exits.
Thus the human stop/go boundary is a completed group, whereas the finer units
are durable recovery boundaries inside that group.

The two values of $G$ were fixed separately for the reported targets from
measured unit throughput and the intended human-scale run duration.  They are
campaign parameters for these two completed computations, not a proposed
general rule for targets that were never run.

The grouping is an operational partition only.  It neither changes the
incremental predicate nor filters the objects being counted.  The $n=21$
groups had recorded durations totaling $8757.72$ seconds.  For $n=22$, the
eight groups totaled $31789.90$ seconds; seven were within the intended
30--90 minute range, and one took $5572.62$ seconds, slightly above 90
minutes.  These times are target-run telemetry reported in the post-result
validation records, not a complexity claim or a forecast for an unexecuted
target~\cite{cortiA369431Package2026}.

\section{Implementation and persistent evidence}
\label{sec:implementation}

The production program is a single C source file compiled in C17 mode,
\path{src/a369431.c}, built by \path{src/Makefile}.  It is serial and contains
the fixed basis, the 20 order-$4$ roots, the order-$5$ frontier construction,
and the recursive counter.  Permutation entries use fixed-size
\texttt{uint8\_t} arrays, while leaf and node counts use
the GCC/Clang extension \texttt{unsigned \_\_int128} with checked addition
and explicit decimal conversion.  Any failed checked addition terminates the
unit with a nonzero outcome and records \texttt{counter-overflow}; silent
wraparound is not accepted.  The source-level campaign cap is $n\leq24$, well
inside both the fixed array capacity and the \texttt{uint8\_t} value range.
No recurrence generator, runtime pattern
parser, external executable, subprocess, or general scheduler participates
in the count.

At each recursive node, the program inserts the next maximum in all positions,
applies the predicate of Section~\ref{sec:mathematics}, and recurses only on
accepted children.  Each call on an accepted permutation, including the
order-$5$ unit root and a target-order permutation, contributes one node;
rejected candidate insertions contribute none.  A target-order node also
contributes one leaf.  Thus, for the reported targets,
$\operatorname{nodes}(n)=\sum_{m=5}^{n}a(m)$.  The node counter supports
telemetry and gives an additional aggregate invariant but does not affect
acceptance.

Each target has one append-only manifest.  Before counting a unit, the program
appends a \texttt{running} row; after successful completion it appends a
\texttt{complete} row containing leaf and node totals, elapsed time, coverage
label, source and binary identities, timestamp, and stop reason.  The manifest
is locked during an invocation, checked on every restart, and refuses rows
with foreign provenance or an invalid state transition.  A completed unit is
never silently overwritten.  When and only when all 66 units are complete,
the program sums their counts and publishes one no-clobber aggregate TSV for
the target.  The aggregate records complete group, unit, and atom coverage as
well as the manifest identity and method.

The manifest schema is the 12-column TSV sequence
\texttt{sequence}, \texttt{target}, \texttt{segment}, \texttt{coverage},
\texttt{status}, \texttt{leaves}, \texttt{nodes}, \texttt{elapsed\_s},
\texttt{source\_sha256}, \texttt{binary\_sha256}, \texttt{epoch}, and
\texttt{stop\_reason}.  For an order-$5$ unit, the \texttt{coverage} value is
the explicit string \texttt{root=$r$,insert5=$p$}, identifying its order-$4$
parent and zero-based insertion position.  The aggregate schema records
sequence, target, coverage, leaf and node totals, source, binary, and manifest
SHA-256 values, method, and recurrence status.  These schemas are data
contracts for the single manifest and aggregate, not separate per-unit files.

Long invocations print an immediate unit-start line, periodic heartbeats, and
a unit-closure line with counts and elapsed time.  This terminal telemetry
does not create auxiliary log trees and is not part of the manifest/aggregate
truth boundary.  The durable public evidence therefore remains compact: one
monolithic source, one Makefile, one manifest and one aggregate per target,
the consolidated term tables, and the validation reports.

The bounded local verification commands are
\begin{verbatim}
make -C src clean test
make -C src sanitize
make -C src clang
make -C src static
\end{verbatim}
from the public package root.  The archived production interface processes a
single group with \texttt{src/a369431 run TARGET GROUP}, or the next incomplete
group with \texttt{src/a369431 run TARGET next}.  Reproduction of a complete
large campaign is intentionally distinct from the short validation commands:
starting from an empty target manifest, four successful invocations of
\texttt{src/a369431 run 21 next} close $n=21$, and eight successful
invocations of \texttt{src/a369431 run 22 next} close $n=22$; the final
invocation publishes the aggregate automatically.  Reissuing \texttt{next}
after an interruption reruns only the incomplete logical unit within the
selected group and preserves completed units.

The target workstation ran Linux Mint 22.3 on an AMD Ryzen 9 7940HS
(8 cores/16 threads) with 64~GB of installed DDR5 RAM.  The reviewed Makefile
default was \texttt{-O3 -std=c17} plus the warning flags printed in
\path{src/Makefile}.  Exact production kernel, compiler version, and any
operator-supplied flag override were not retained; therefore the elapsed times
in Section~\ref{subsec:groups} are provenance telemetry, not portable
performance benchmarks.  This missing metadata is a reproducibility
limitation rather than a value-correction claim.

\section{Validation and certification boundary}
\label{sec:validation}

Validation addressed the mathematical predicate, implementation, persistence,
and final arithmetic through several partially independent checks.  The main
checks and their recorded scopes are summarized in Table~\ref{tab:validation}.

For each target manifest, the post-result audit found 66 \texttt{running} and
66 \texttt{complete} rows, no missing or duplicate complete unit, and exact
agreement between the sum of completed rows and the final aggregate.  The
source and binary identities in every row matched those of the reviewed
canonical production candidate.  At campaign closure, the then-current b-file
and consolidated result table agreed term by term through $n=22$.  The public
b-file has subsequently been extended by the separate formula-derived process
in Section~\ref{sec:recurrence}; the certified-term and exhaustive-result TSVs
deliberately remain bounded at $n=22$~\cite{cortiA369431Package2026}.

The small direct predicate is mathematically different from the incremental
test but is included in the same production source for selftest purposes.
An external Claude review additionally reported an LLM-generated Python brute
force comparison through $n=9$, independent of the C implementation but not
an independent human-authored verification.  Neither check is a second
exhaustive enumeration at $n=21$ or $n=22$.  Likewise, the LLM reviews are
code and artifact reviews, not independent mathematical peer review and not
a second production campaign.

Accordingly, ``validated'' or ``certified'' in this paper has an explicit
artifact-bounded meaning: Proposition~\ref{prop:insertion} supplies the
mathematical coverage argument; the reviewed implementation supplies the
enumerator; and the manifests and aggregates certify completion and arithmetic
consistency for that implementation.  The residual common-mode risk is an
implementation error not detected by the direct small-order oracle, known-term
replay, compiler/sanitizer/static checks, or external code reviews.

The sanitizer stage included AddressSanitizer, the compiler-instrumentation
method of Serebryany, Bruening, Potapenko, and Vyukov, together with
UndefinedBehaviorSanitizer~\cite{serebryany2012}.  These tools test important
memory-safety and undefined-behavior boundaries, but their successful execution
is not by itself a proof of algorithmic correctness.

The theorem-critical source, binary, manifest, and aggregate SHA-256
identities are recorded in the package's validation ledger and post-result
reports~\cite{cortiA369431Package2026}.  Keeping the inventory in those
machine-audit documents avoids turning the scientific narrative into a hash
catalogue.  Byte identities complement, but do not replace, a versioned
public release and archival DOI.

\clearpage
\begin{table}[!t]
\centering
\caption{Validation evidence recorded for the canonical production source.}
\label{tab:validation}
\begin{tabular}{L{0.36\textwidth}L{0.47\textwidth}}
\toprule
Check & Recorded scope and result \\
\midrule
Root and frontier construction & The explicit 20 roots match the direct
  order-$4$ oracle; the valid order-$5$ frontier has 66 units. \\
Incremental predicate comparison & Exhaustive equality with the direct
  four-tuple predicate through order $9$. \\
Direct small-order oracle & Known values reproduced through order $9$. \\
Maximum-insertion replay & Known values reproduced through order $12$;
  external review also recorded grouped known-term checks at orders $12$ and
  $16$. \\
Node-count invariant & Summing the consolidated terms according to
  $\operatorname{nodes}(n)=\sum_{m=5}^{n}a(m)$ reproduced both aggregate node
  totals. \\
Persistence selftests & Manifest parsing, state transitions, interrupted
  groups, no-clobber aggregate behavior, group coverage, and numeric parsing
  passed bounded fixtures. \\
Compiler and analysis checks & GCC and Clang builds/selftests, ASan/UBSan,
  \texttt{clang-tidy}, and \texttt{cppcheck} passed under the recorded
  configurations~\cite{llvmSanitizers,llvmClangTidy,cppcheckManual}. \\
Campaign closure & For each of $n=21$ and $n=22$, all 66 units were complete
  exactly once; manifest sums equaled aggregate leaf and node totals; coverage
  was $20/20$ atoms. \\
External LLM-assisted reviews & ChatGPT and Claude reports examined the frozen
  grouping/publication candidate and recorded no reproducible defect in their
  stated scopes; this is review history, not independent peer review or
  numerical evidence. \\
\bottomrule
\end{tabular}
\end{table}
\FloatBarrier

\section{Computational results}
\label{sec:results}

The final aggregate records are shown in Table~\ref{tab:results}.  Here
``nodes'' is the inclusive accepted-node count defined in
Section~\ref{sec:implementation}; it is a diagnostic total distinct from the
number of target-order leaves.

\begin{table}[htbp]
\centering
\caption{Canonical aggregate results.}
\label{tab:results}
\begin{tabular}{rrrrr}
\toprule
$n$ & Groups & Units & Leaves $a(n)$ & Nodes \\
\midrule
21 & 4 & 66 & $8321204878$ & $12094019604$ \\
22 & 8 & 66 & $26674199972$ & $38768219576$ \\
\bottomrule
\end{tabular}
\end{table}

\begin{theorem}[Artifact-bounded computational result]
\label{thm:main}
Subject to the correct execution of the reviewed production implementation
described in Section~\ref{sec:implementation} and the completed canonical
manifests described in Section~\ref{sec:validation}, exhaustive
maximum-insertion enumeration gives
\[
 \lvert\operatorname{Av}_{21}(B)\rvert=8321204878
\]
and
\[
 \lvert\operatorname{Av}_{22}(B)\rvert=26674199972.
\]
\end{theorem}

\begin{proof}[Computational proof]
Proposition~\ref{prop:insertion} proves that the maximum-insertion tree accepts
exactly the avoiders and that every target permutation has a unique path.  The
20 order-$4$ atoms cover all possible paths, and their 66 order-$5$ children
give a disjoint refinement.  For each target, the canonical manifest contains
one completed count for every one of these 66 units.  Post-result validation
found no missing or duplicated complete unit and verified that the sums of
the unit leaf and node counts equal the corresponding aggregate row.  The
aggregate leaf totals are those displayed in Table~\ref{tab:results}, which
establishes the two finite computational claims within the certification boundary
of Section~\ref{sec:validation}.
\end{proof}

The theorem gives a computation independent of the structural result below.
The package's \path{data/certified_terms.tsv} and
\path{results/a369431_terms.tsv} stop at $n=22$ so that the exhaustive
certification boundary remains visible.

\section{External structural result and formula-derived b-file}
\label{sec:recurrence}

The PermPAL record for this exact class reports the rational generating
function
\begin{equation}
 A(x)=\sum_{n\geq0}a(n)x^n
 =\frac{(x-1)(x^2-3x+1)}{4x^3-7x^2+5x-1},
 \label{eq:permpal-gf}
\end{equation}
together with a machine-generated proof tree and counting sequence
\cite{albertEtAl2026,permpalA369431Class}.  The record states that its
\emph{Point Placements} specification has 51 rules and was found by the
PermPAL system on January 18, 2022.  This source was located only after
the $n=21$ and $n=22$ campaigns and the initial manuscript had been completed.
It was therefore absent from the production design and canonical execution
path.  The paper does not claim to have discovered or reproved
\eqref{eq:permpal-gf}.

Multiplying numerator and denominator by $-1$ rewrites
\eqref{eq:permpal-gf} as
\[
 (1-5x+7x^2-4x^3)A(x)=1-4x+4x^2-x^3.
\]
Coefficient comparison gives $a(0),a(1),a(2),a(3)=1,1,2,6$ and, for
$n\geq4$,
\begin{equation}
 a(n)=5a(n-1)-7a(n-2)+4a(n-3).
 \label{eq:recurrence}
\end{equation}
At the two independently enumerated endpoints, the formula gives the
post-hoc equalities
\begin{align*}
5(2595858574)-7(809796400)+4(252621702) &= 8321204878,\\
5(8321204878)-7(2595858574)+4(809796400) &= 26674199972.
\end{align*}
This agreement is not a generator, pruning rule, acceptance test, independent
campaign certificate, or proof of Theorem~\ref{thm:main}.

For the separate data-extension task, however, the prior PermPAL formula is
the declared generator.  The Python 3.12.3 tool
\path{tools/generate_b369431.py} iterates \eqref{eq:recurrence} with exact
integers through $n=400$, starting from the four displayed values
$a(0),\ldots,a(3)=1,1,2,6$.  Independently, PARI/GP 2.15.4 expands
\eqref{eq:permpal-gf} directly as a formal power series; its script does not
iterate the Python implementation~\cite{pythonDocs,pariGp}.  The two outputs
agree byte for byte for all 401 coefficients.  They also agree with the 101
coefficients displayed by PermPAL for $0\leq n\leq100$ and with the previous
local b-file for $0\leq n\leq22$.

Thus the recurrence reproduces every term in the pre-campaign OEIS range
$0\leq n\leq20$ as well as the two C endpoints.  In this paper,
``formula-derived'' means a coefficient of the cited rational series,
calculated and cross-checked locally; its validity remains conditional on that
external generating function and does not carry the exhaustive C certificate.

The cutoff at $n=400$ is editorial, not mathematical.  Since the denominator
of \eqref{eq:permpal-gf} has nonzero constant term, the rational function has
a unique formal power-series expansion at the origin, and
\eqref{eq:recurrence} determines every later coefficient from three preceding
ones.  Thus, conditional on the correctness of the cited PermPAL generating
function, the calculation can be continued to any prescribed finite index;
only integer size, running time, memory, and output size impose practical
limits.  We stop at $400$ to keep the public data proportionate to the paper.

Representative new b-file rows are
\[
\begin{array}{c@{\qquad}r}
n & a(n)\\
\hline
23 & 85506000010\\
24 & 274095419758\\
25 & 878631898608\\
40 & 34070992412092462900\\
50 & 3903342517082139278233270
\end{array}
\]
The b-file continues through $a(400)$, a 202-digit integer.  Its full value
belongs in the machine-readable data rather than the running text.  The
validation report records the commands, comparisons, file format checks, and
SHA-256 identities.  This is a reproducible execution and verification of a
previously published formula, not a claim of priority for the resulting
coefficients.

\section{Limitations}
\label{sec:limitations}

The reported certification is finite and tied to the archived source,
manifests, aggregate tables, and validation records.  There is no second
independent exhaustive implementation reproducing the complete $n=21$ and
$n=22$ searches.  The direct oracle and LLM-generated Python check cover only
small orders, while the known-term and grouped replays cover bounded larger
orders.  External LLM-assisted reviews strengthen the source and artifact
audit but do not replace conventional mathematical peer review.

The exact production kernel, compiler version, and effective operator build
flags were not retained, so the reported elapsed times cannot be reproduced as
benchmarks even though source, binary identity, data, and completion records
are pinned.

The exhaustive campaign ends at $a(22)$; no maximum-insertion campaign for
$a(23)$ or any higher term was executed.  Rows $23..400$ of the b-file have
the weaker and different status ``derived from the PermPAL generating
function.''  Agreement between Python and PARI/GP protects against elementary
transcription or implementation errors but is not an independent proof of the
underlying 51-rule PermPAL specification cited in
Section~\ref{sec:recurrence}.  Possible future work includes
a human-readable derivation or independent audit of that specification and an
independent exhaustive reproduction of the two finite counts; neither has
been performed here.

\section{Conclusion}
\label{sec:conclusion}

Formula-free maximum insertion independently verifies
\[
 a(21)=8321204878,
 \qquad
 a(22)=26674199972.
\]
The key computational structure is the separation between mathematical
coverage and operational recovery: 20 order-$4$ atoms prove coverage, 66
order-$5$ units provide durable disjoint subtasks, and only a few balanced
groups define the operator-facing runs.  This yields a compact auditable
record without changing the enumeration.  The previously published PermPAL
generating function then permits a negligible-cost exact calculation of a
b-file through $n=400$.  Two arbitrary-precision routes and the 101 displayed
PermPAL coefficients cross-check that calculation.  The scientific roles are
therefore deliberately distinct: independent exhaustive evidence at $n=21,22$
and attributed formula-derived data at $n=23,\ldots,400$.

\section*{Data availability}
\begin{sloppypar}

The public package contains the C code, b-file candidate,
formula-generation scripts, certified-term table, result tables, append-only
manifests, aggregate records, validation reports, and the manuscript as both
\path{paper/A369431.tex} and \path{paper/A369431.pdf}.  It is available from
the GitHub repository \url{https://github.com/carcorti/A369431} and the Zenodo
archive \url{https://doi.org/10.5281/zenodo.21485629}.  The formula-derived
extension is reproduced by
\path{tools/generate_b369431.py} and \path{tools/generate_b369431.gp} and
audited in
\path{validation/A369431_formula_derived_bfile_validation.md}.  The exhaustive
evidence includes \path{data/b369431.txt}, the two manifests
\path{validation/a369431_n21_manifest.tsv} and
\path{validation/a369431_n22_manifest.tsv}, and the two aggregate records
\path{results/a369431_n21.aggregate.tsv} and
\path{results/a369431_n22.aggregate.tsv}.  No per-unit directory or external
dataset is needed to audit the reported finite computation.
\end{sloppypar}

\section*{AI use disclosure}

Large language models were used as research, review, and drafting assistants
during the A369431 workflow: OpenAI Codex and ChatGPT, Claude, DeepSeek,
Gemini, Grok, Kimi, Qwen, MiniMax, and Z.  Round-01 manuscript reports also
included Lumo and Meta model families.  The dated reports are retained in the
local prepublication review archive.  Names are given at family or tool level
because exact model-version metadata were not consistently retained in the
earlier local artifacts; the reports preserve the identifiers declared by
each reviewer.

\noindent\textbf{Mathematical analysis.}
DeepSeek contributed preliminary triage, and Codex assisted with reconciling
the sequence definition, the maximum-insertion argument, and the distinction
between the initially empirical recurrence and the later-located PermPAL
result.  Mathematical
claims retained in this paper were checked against the local OEIS snapshot,
the explicit four-case incremental predicate, the b-file, and the aggregate
result tables.  The recurrence was deliberately excluded from the exhaustive
C generation and certification.  It was subsequently used, with attribution,
only for the formula-derived b-file described in
Section~\ref{sec:recurrence}.

\noindent\textbf{Algorithm and code.}
Codex assisted with computational-strategy analysis, validation planning,
code and artifact review, and integration of review evidence.  ChatGPT and
Claude performed external LLM-assisted reviews of the frozen production
candidate; Claude also reported an independent Python small-order brute-force
check, independent of the C implementation but not human-authored.  Acceptance
of the final implementation was tied to the canonical C17
source and Makefile, selftests, direct-oracle and known-term comparisons,
manifest/aggregate equality, compiler and sanitizer runs, static analysis,
and the post-result validation reports.  No conversational LLM assertion is
used as numerical evidence.

\noindent\textbf{Manuscript.}
Codex assisted with drafting, internal adversarial review, comparison of the
11 first-round reports, targeted literature research, successive manuscript
revisions, the author-directed formula-extension revision, the
formula-generation and validation workflow, and comparison of the 11
subsequent review reports.
The reviewers received overlapping manuscript-centered packets and did not
constitute independent human peer review.  Every headline value and coverage statement was checked against
the b-file candidate, certified-term and result TSVs, canonical manifests,
aggregate TSVs, validation summaries, and code-review integration record.
Formula-derived rows were additionally checked by exact Python and PARI/GP
calculations and against PermPAL's 101 displayed coefficients.  The numerical
claims rely on those artifacts and checks, not on generated prose.

The author takes full scientific responsibility for the mathematical
statements, algorithm, implementation, numerical results, validation
boundaries, and conclusions of this paper.

\clearpage
\begin{thebibliography}{99}
\sloppy
\raggedright

\bibitem{albertLintonRuskuc2005}
M. H. Albert, S. Linton, and N. Ru\v{s}kuc,
\emph{The insertion encoding of permutations},
Electronic Journal of Combinatorics \textbf{12} (2005), Article R47.
\url{https://doi.org/10.37236/1944}.

\bibitem{albertEtAl2026}
M. H. Albert, C. Bean, A. Claesson, \'E. Nadeau, J. Pantone, and
H. Ulfarsson,
\emph{Combinatorial Exploration: An Algorithmic Framework for Enumeration},
Memoirs of the American Mathematical Society \textbf{317} (2026), no.~1611.
ISBN 978-1-4704-8000-4.

\bibitem{bona2022}
M. B\'ona,
\emph{Combinatorics of Permutations}, 3rd ed.,
Discrete Mathematics and Its Applications, CRC Press, 2022.
ISBN 978-0-367-22258-1.
\url{https://www.routledge.com/Combinatorics-of-Permutations/Bona/p/book/9780367222581}.

\bibitem{cortiA369431Package2026}
C. Corti,
\emph{OEIS A369431: Independent Maximum-Insertion Enumeration through
$n=22$ and Formula-Derived Data through $n=400$},
GitHub and Zenodo release package, 2026.
\url{https://github.com/carcorti/A369431}.
\url{https://doi.org/10.5281/zenodo.21485629}.

\bibitem{cppcheckManual}
Cppcheck project,
\emph{Cppcheck manual}, official documentation.
\url{https://cppcheck.sourceforge.io/manual.html}.
Accessed July 21, 2026.

\bibitem{kitaev2011}
S. Kitaev,
\emph{Patterns in Permutations and Words},
Monographs in Theoretical Computer Science, Springer, Berlin and Heidelberg,
2011.
\url{https://doi.org/10.1007/978-3-642-17333-2}.

\bibitem{llvmClangTidy}
LLVM Project,
\emph{clang-tidy---Clang extra tools documentation}.
\url{https://clang.llvm.org/extra/clang-tidy/}.
Accessed July 21, 2026.

\bibitem{llvmSanitizers}
LLVM Project,
\emph{UndefinedBehaviorSanitizer documentation}.
\url{https://clang.llvm.org/docs/UndefinedBehaviorSanitizer.html}.
Accessed July 21, 2026.

\bibitem{oeisA369431}
OEIS Foundation Inc.,
\emph{Sequence A369431: Number of permutations of $[n]$ avoiding 1234, 1324,
1342, and 2413},
The On-Line Encyclopedia of Integer Sequences,
\url{https://oeis.org/A369431}.
Submitted by Matt Slattery-Holmes, January 23, 2024; terms $a(13)$--$a(20)$
contributed by Martin Ehrenstein, February 24, 2024.  Local project snapshot
consulted July 21, 2026.

\bibitem{permpalA369431Class}
Permutation Pattern Avoidance Library (PermPAL),
\emph{$\operatorname{Av}(1234,1324,1342,2413)$}, record 14688.
\url{https://www.permpal.com/perms/basis/0123_0213_0231_1302/}.
The page reports a 51-rule \emph{Point Placements} specification found by the
PermPAL system on January 18, 2022; accessed July 21, 2026.

\bibitem{pariGp}
The PARI Group,
\emph{PARI/GP User's Guide, version 2.15.4},
Universit\'e de Bordeaux.
\url{https://pari.math.u-bordeaux.fr/pub/pari/manuals/2.15.4/users.pdf}.

\bibitem{pythonDocs}
Python Software Foundation,
\emph{Python 3.12 documentation}.
\url{https://docs.python.org/3.12/}.
Accessed July 21, 2026.

\bibitem{serebryany2012}
K. Serebryany, D. Bruening, A. Potapenko, and D. Vyukov,
\emph{AddressSanitizer: A fast address sanity checker},
in \emph{2012 USENIX Annual Technical Conference (USENIX ATC 12)},
USENIX Association, 2012, pp.~309--318.
\url{https://www.usenix.org/conference/atc12/technical-sessions/presentation/serebryany}.

\bibitem{simionSchmidt1985}
R. Simion and F. W. Schmidt,
\emph{Restricted permutations},
European Journal of Combinatorics \textbf{6} (1985), no.~4, 383--406.
\url{https://doi.org/10.1016/S0195-6698(85)80052-4}.

\bibitem{weissteinPermutationPattern}
E. W. Weisstein,
\emph{Permutation Pattern},
MathWorld---A Wolfram Resource.
\url{https://mathworld.wolfram.com/PermutationPattern.html}.
Accessed July 21, 2026.

\end{thebibliography}

\end{document}
